% 底本 EGAIV.1.pdf，上传PDF第135--139页，原分册印刷第136--140页。
\subsection{非完美性模与特征同态}
\source{第四卷第一部分，分册印刷第136--147页；上传PDF第135--146页，20.6全文}
\egas{20.6.1}{定义}
给定两个环同态 $u:A\to B$、$v:B\to C$，称 $C$-模
\[
 \Upsilon_{C/B/A}
 =\Ker\bigl(v_{C/B/A}:\Omega^1_{B/A}\otimes_B C\longrightarrow\Omega^1_{C/A}\bigr)
\]
为\emph{$B$-代数 $C$ 相对于 $A$ 的非完美性模}。因而根据定义及（20.5.7），有正合列
\begin{equation}
\begin{aligned}
0\longrightarrow\Upsilon_{C/B/A}
 &\longrightarrow\Omega^1_{B/A}\otimes_B C
 \xrightarrow{v_{C/B/A}}\Omega^1_{C/A}\\
 &\xrightarrow{u_{C/B/A}}\Omega^1_{C/B}\longrightarrow0.
\end{aligned}
\tag{20.6.1.1}
\end{equation}
当 $A=\mathbf Z$ 时，$\Omega^1_{B/A}$、$\Omega^1_{C/A}$ 分别是绝对微分模
$\Omega^1_B$、$\Omega^1_C$，此时将 $\Upsilon_{C/B/\mathbf Z}$ 简记为
$\Upsilon_{C/B}$。若 $B,C$ 均为某个素域 $P$ 上的代数，则
$\Upsilon_{C/B/P}=\Upsilon_{C/B}$。

设 $R,S$ 分别为 $B,C$ 的乘法子集，且 $R$ 的像包含在 $S$ 中。由正合列
（20.6.1.1）及（20.5.9）得到典范识别
\begin{equation}
 \Upsilon_{S^{-1}C/R^{-1}B/A}
 \simeq S^{-1}\Upsilon_{C/B/A}.
 \tag{20.6.1.2}
\end{equation}

\egas{20.6.2}{命题}
若 $C$ 是相对于 $A$ 形式光滑的 $B$-代数，特别地，若 $C$ 是形式光滑的
$B$-代数，则 $\Upsilon_{C/B/A}=0$。

这由（20.5.7，ii）得到。
\cmt{J10-N01}{这里只断言形式光滑性蕴涵非完美性模消失，未断言对一般环的逆命题。
例如 $A=B=k$、$C=k[\epsilon]/(\epsilon^2)$ 时该非完美性模为零，
但 $C$ 对 $k$ 并不形式光滑。域扩张的等价判据是下一条在特定条件下的结论。}

\egas{20.6.3}{命题}
设 $k$ 为域，$K$ 为 $k$ 的扩张。$K/k$ 可分，当且仅当 $\Upsilon_{K/k}=0$。
换言之，当且仅当典范映射
\[
 \Omega^1_k\otimes_k K\longrightarrow\Omega^1_K
\]
单射；也等价于每个从 $k$ 到 $K$ 的导子都能延拓为 $K$ 到自身的导子
（20.4.8）。

由（19.6.1），$K/k$ 可分等价于 $K$ 是形式光滑的 $k$-代数。
设 $P$ 为 $k$ 的素域。$K/P$ 可分，因此 $K$ 对 $P$ 形式光滑；
于是 $K$ 对 $k$ 形式光滑，等价于 $K$ 对 $k$ 相对于 $P$ 形式光滑。
最后，根据（20.5.7，ii），后一条件等价于
\[
 v_{K/k/P}:\Omega^1_{k/P}\otimes_k K\longrightarrow\Omega^1_{K/P}
\]
有左逆。由于 $K$ 是域，这又等价于其核 $\Upsilon_{K/k}$ 为零。
导子形式的表述则来自（20.4.8）及向量空间中单射的分裂性。

\egap{20.6.4}
考虑交换的环同态图
\begin{equation}
\begin{tikzcd}
 A'\arrow[r,"u'"]&B'\arrow[r,"v'"]&C'\\
 A\arrow[u]\arrow[r,"u"']&B\arrow[u]\arrow[r,"v"']&C.\arrow[u]
\end{tikzcd}
\tag{20.6.4.1}
\end{equation}
相应图（20.5.7.3）的交换性给出一个唯一的 $C$-线性同态
\begin{equation}
 \Upsilon_{C/B/A}\longrightarrow\Upsilon_{C'/B'/A'},
 \tag{20.6.4.2}
\end{equation}
其中右端经 $C\to C'$ 视为 $C$-模；该同态使下图交换且两行正合：
\begin{equation}
\begin{gathered}
\begin{tikzcd}[column sep=small]
0\arrow[r]&\Upsilon_{C/B/A}\arrow[r]\arrow[d]&
 \Omega^1_{B/A}\otimes_B C\arrow[r,"v_{C/B/A}"]\arrow[d]&
 \Omega^1_{C/A}\arrow[d]\\
0\arrow[r]&\Upsilon_{C'/B'/A'}\arrow[r]&
 \Omega^1_{B'/A'}\otimes_{B'} C'\arrow[r,"v'_{C'/B'/A'}"']&
 \Omega^1_{C'/A'}
\end{tikzcd}\\[-.2em]
\begin{tikzcd}[column sep=large]
\Omega^1_{C/A}\arrow[r,"u_{C/B/A}"]\arrow[d]&\Omega^1_{C/B}\arrow[r]\arrow[d]&0\\
\Omega^1_{C'/A'}\arrow[r,"u'_{C'/B'/A'}"']&\Omega^1_{C'/B'}\arrow[r]&0 .
\end{tikzcd}
\end{gathered}
\tag{20.6.4.3}
\end{equation}
这里为排版把同一长图分成两段，两段共有的 $\Omega^1_{C/A}$ 及其竖箭头相同。
给出（20.6.4.2），也等价于给出一个 $C'$-线性同态
\begin{equation}
 \Upsilon_{C/B/A}\otimes_C C'\longrightarrow\Upsilon_{C'/B'/A'}.
 \tag{20.6.4.4}
\end{equation}
它与典范映射 $\Upsilon_{C/B/A}\to\Upsilon_{C/B/A}\otimes_C C'$ 复合后，
就是（20.6.4.2）。这些同态满足传递性，因而
$\Upsilon_{C/B/A}$ 是关于三元组 $(A,B,C)$ 的函子。

\egap{20.6.5}
在（20.6.1）的条件下，为便于后文，引入 $C$-模的链复形 $K_\bullet(C/B/A)$。
除次数 $0,1$ 外，其各项均为零，并令
\begin{equation}
 \begin{cases}
 K_0(C/B/A)=\Omega^1_{C/A},\\
 K_1(C/B/A)=\Omega^1_{B/A}\otimes_B C .
 \end{cases}
 \tag{20.6.5.1}
\end{equation}
微分 $K_1\to K_0$ 取为 $v_{C/B/A}$。于是有典范识别
\begin{equation}
 H_0(K_\bullet(C/B/A))=\Omega^1_{C/B},\qquad
 H_1(K_\bullet(C/B/A))=\Upsilon_{C/B/A}.
 \tag{20.6.5.2}
\end{equation}

\egas{20.6.6}{命题}
在（20.6.1）的假设下，对每个 $C$-模 $L$ 有典范的 $C$-线性同构
\begin{equation}
 H^0(K_\bullet(C/B/A),L)\simeq\Der_B(C,L),
 \tag{20.6.6.1}
\end{equation}
\begin{equation}
 H^1(K_\bullet(C/B/A),L)\simeq\Exalcom_{B/A}(C,L).
 \tag{20.6.6.2}
\end{equation}
事实上，由（20.4.8）及（20.5.6），上链复形
$\Hom_C(K_\bullet(C/B/A),L)$ 就是
\[
 \cdots\longrightarrow0\longrightarrow\Der_A(C,L)
 \xrightarrow{v^0}\Der_A(B,L)\longrightarrow0\longrightarrow\cdots,
\]
其中 $v^0$ 为（20.2.1）中的限制映射。命题因而来自正合列（20.2.4.2）
以及上同调的定义
\begin{equation}
 H^i(K_\bullet,L)=H^i\bigl(\Hom_C(K_\bullet,L)\bigr).
 \tag{20.6.6.3}
\end{equation}
\cmt{J10-N02}{这里的 $H^i(K_\bullet,L)$ 是所写两项 $\Hom$ 复形的普通上同调。
正文不把它替换为导出函子 $\Ext$，也不假设 $K_0,K_1$ 已是投射模。}

若有交换的环同态图
\[
\begin{tikzcd}
 A'\arrow[r]&B'\arrow[r]&C'\\
 A\arrow[u]\arrow[r]&B\arrow[u]\arrow[r]&C,\arrow[u]
\end{tikzcd}
\]
则（20.5.4.3）的双同态定义复形的双同态
\begin{equation}
 K_\bullet(C/B/A)\longrightarrow K_\bullet(C'/B'/A').
 \tag{20.6.6.4}
\end{equation}
由此在同调或上同调上诱导的双同态，经（20.6.5.2）、（20.6.6.1）、
（20.6.6.2）的识别，分别就是此前在（20.5.4.3）、（20.6.4.2）、
（20.2.1）和（18.3.7.2）中定义的同态。最后一项须同时使用正合列
（20.2.4.2）中算子 $\partial$ 的定义。

\egap{20.6.7}
对一个 $C$-模复形 $K_\bullet$ 和一个 $C$-模 $L$，存在典范同态
\[
 \alpha_i:H^i(K_\bullet,L)\longrightarrow\Hom_C(H_i(K_\bullet),L)
 \qquad\text{（M，第四章，§6）}.
\]
在这里，$\alpha_1$ 可直接由限制映射
\[
 \Hom_C(K_1(C/B/A),L)\longrightarrow
 \Hom_C(H_1(K_\bullet(C/B/A)),L)
\]
对 $\Hom_C(K_0(C/B/A),L)$ 的像取商而得到；因为
$H_1(K_\bullet(C/B/A))=\Ker(K_1\to K_0)$，该像的限制为零。
结合（20.6.6.2）与（20.6.5.2），得到典范的 $C$-线性同态
\begin{equation}
 \Exalcom_{B/A}(C,L)\longrightarrow\Hom_C(\Upsilon_{C/B/A},L).
 \tag{20.6.7.1}
\end{equation}

它可以具体描述如下。由（20.2.4.2），每个以 $L$ 为核且为 $A$-平凡的
$C$ 的 $B$-扩张，均由某个 $A$-导子 $D:B\to L$ 给出。
根据（20.4.8），它对应于一个 $C$-线性同态
$f:\Omega^1_{B/A}\otimes_B C\to L$。将该扩张的类映到
$f|_{\Upsilon_{C/B/A}}$；这个限制只依赖扩张类，不依赖于 $D$ 的选择。

\egas{20.6.8}{定义}
设 $u:A\to B$、$v:B\to C$ 为两个环同态。对于以 $C$-模 $L$ 为核、
且为 $A$-平凡的 $C$ 的任一 $B$-扩张 $E$（18.3.7），称其类在
（20.6.7.1）下的像为\emph{$E$ 的特征同态}，记作
\[
 \chi_E:\Upsilon_{C/B/A}\longrightarrow L .
\]
还可以用以下方式定义这个同态。

\egas{20.6.9}{命题}
设 $E$ 为以 $C$-模 $L$ 为核、且为 $A$-平凡的 $C$ 的一个 $B$-扩张。
则下图交换且两行正合：
\begin{equation}
\begin{tikzcd}[column sep=small]
0\arrow[r]&\Upsilon_{C/B/A}\arrow[r]\arrow[d,"\chi_E"']&
 \Omega^1_{B/A}\otimes_B C\arrow[r,"v_{C/B/A}"]\arrow[d,"q_{E/B/A}\otimes1_C"']&
 \Omega^1_{C/A}\arrow[d,"1"]&\\
0\arrow[r]&L\arrow[r,"\delta_{C/E/A}"']&
 \Omega^1_{E/A}\otimes_E C\arrow[r,"p_{C/E/A}"']&
 \Omega^1_{C/A}\arrow[r]&0 .
\end{tikzcd}
\tag{20.6.9.1}
\end{equation}
这里 $q:B\to E$ 给出 $E$ 的 $B$-扩张结构，$p:E\to C$ 为增广同态。

下行是与 $A\to E$、$p:E\to C=E/L$ 相应的（20.5.12.5）。
由于 $p$ 满射且扩张为 $A$-平凡，由（20.5.12，iii），该行正合且分裂。
右方方块的交换性来自 $v=pq$ 及（20.5.2.7）。因此
$q_{E/B/A}\otimes1_C$ 将 $\Upsilon_{C/B/A}$ 映入
$\Ker(p_{C/E/A})$，后者经 $\delta_{C/E/A}$ 识别为 $L$。

为核对这个限制正是 $\chi_E$，取 $p$ 的一个 $A$-代数截面 $h:C\to E$，
并记核的嵌入为 $j:L\to E$。于是
\[
 q(b)=h(v(b))+j(D(b))\qquad(b\in B),
\]
其中 $D:B\to L$ 为定义这个 $B$-扩张的 $A$-导子。
写成 $D=f_0d_{B/A}$，其中 $f_0:\Omega^1_{B/A}\to L$ 为 $B$-线性同态。
对
\[
 z=\sum_i d_{B/A}(b_i)\otimes c_i
 \in\Omega^1_{B/A}\otimes_B C
\]
有
\[
\begin{aligned}
(q_{E/B/A}\otimes1_C)(z)
={}&\sum_i d_{E/A}(h(v(b_i)))\otimes c_i\\
 &+\sum_i d_{E/A}(j(f_0(d_{B/A}(b_i))))\otimes c_i .
\end{aligned}
\]
第一项为 $(h_{E/C/A}\otimes1_C)(v_{C/B/A}(z))$。
当 $z\in\Upsilon_{C/B/A}$ 时它为零；第二项按定义为
$\delta_{C/E/A}(f(z))$，其中 $f$ 是 $f_0$ 对应的 $C$-线性同态。
而 $f(z)=\chi_E(z)$，这证明左方方块交换。

若只假定 $\delta_{C/E/A}$ 单射，而不要求它有左逆，这个解释仍可用于定义
$\chi_E$：先限制 $q_{E/B/A}\otimes1_C$ 到 $\Upsilon_{C/B/A}$，再由
$\delta_{C/E/A}$ 将其像识别为 $L$。

\egas{20.6.10}{推论}
设 $\mathfrak k$ 为 $B$ 的理想，$C=B/\mathfrak k$、$E=B/\mathfrak k^2$。
若 $E$ 是以 $\mathfrak k/\mathfrak k^2$ 为核的 $C$ 的一个 $A$-平凡的
$B$-扩张，则特征同态
\[
 \chi_E:\Upsilon_{C/B/A}\longrightarrow\mathfrak k/\mathfrak k^2
\]
双射。

事实上，在图（20.6.9.1）中，右边两个竖箭头均双射（20.5.12，ii），
于是它们在核上诱导的同态也双射。

\egas{20.6.11}{定理}
设 $u:A\to B$、$v:B\to C$ 为两个环同态，$L$ 为 $C$-模。
假定满足下列条件之一：
\begin{enumerate}[label=(\roman*)]
\item $L$ 为内射 $C$-模。
\item $\Upsilon_{C/B/A}$ 是 $\Omega^1_{B/A}\otimes_B C$ 的直和因子，
且 $u_{C/B/A}:\Omega^1_{C/A}\to\Omega^1_{C/B}$ 有右逆。
\end{enumerate}
则典范同态
\[
 \Exalcom_{B/A}(C,L)\longrightarrow\Hom_C(\Upsilon_{C/B/A},L)
\]
双射。特别地，若 $\Omega^1_{C/A}$、$\Omega^1_{C/B}$ 都是投射 $C$-模，
则对每个 $L$，该同态双射。

为把 $\alpha_1$ 的定义直接用于这两个情形，记
\[
\begin{gathered}
 M=\Omega^1_{B/A}\otimes_B C,\qquad N=\Omega^1_{C/A},\\
 Y=\Ker(M\to N),\qquad I=\Im(M\to N),\qquad Q=N/I .
\end{gathered}
\]
于是
\[
 H^1(K_\bullet,L)
 =\Hom_C(M,L)\big/\Im\bigl(\Hom_C(N,L)\bigr),
\]
而 $\alpha_1$ 由限制到 $Y$ 给出。若 $L$ 内射，每个 $Y\to L$ 都延拓到
$M$，故 $\alpha_1$ 满射；一个 $M\to L$ 若在 $Y$ 上为零，便分解为
$M\to I\to L$，且 $I\to L$ 可延拓到 $N$，故 $\alpha_1$ 也单射。
若满足 $(ii)$，则 $Y\subseteq M$ 分裂且 $I\subseteq N$ 分裂，
同样的两次延拓仍成立。

还应指出，条件 $(ii)$ 恰好是对所有 $C$-模 $L$ 都使 $\alpha_1$ 双射的
充要条件：取 $L=Y$，满射性给出 $M\to Y$ 对包含的左逆；再取 $L=I$，
单射性使 $M\to I$ 延拓为 $N\to I$，后者是包含 $I\to N$ 的左逆。
于是 $N\to Q$ 有右逆。\cmt{J10-G01}{原文将这部分归结为
$\alpha_1$ 的定义及Bourbaki的分裂判据。正文展开了两次延拓的论证，
但没有改变条件，也没有把“对给定 $L$ 双射”误写为足以推出两个分裂条件；
必要性明确要求对所有 $L$ 双射。}

若 $\Omega^1_{C/A}$ 与 $\Omega^1_{C/B}$ 都投射，则
\[
 0\longrightarrow I\longrightarrow\Omega^1_{C/A}
 \longrightarrow\Omega^1_{C/B}\longrightarrow0
\]
分裂，故 $I$ 是投射模；进而
\[
 0\longrightarrow\Upsilon_{C/B/A}\longrightarrow
 \Omega^1_{B/A}\otimes_B C\longrightarrow I\longrightarrow0
\]
也分裂，从而满足 $(ii)$。

\egas{20.6.12}{推论}
若 $C$ 是形式光滑的 $A$-代数，则有典范同态
\begin{equation}
 \Exalcom_B(C,L)\longrightarrow\Hom_C(\Upsilon_{C/B/A},L).
 \tag{20.6.12.1}
\end{equation}
当（20.6.11）的 $(i)$ 或 $(ii)$ 成立时，它双射。

由 $C$ 的假设，
$\Exalcom_B(C,L)=\Exalcom_{B/A}(C,L)$；
（20.6.12.1）就是（20.6.7.1）。

\egas{20.6.13}{推论}
若 $C$ 为形式光滑的 $A$-代数，且 $\Omega^1_{C/B}$ 是投射 $C$-模，
则（20.6.12.1）双射。

因为在离散拓扑下，$C$ 对 $A$ 形式光滑蕴涵 $\Omega^1_{C/A}$ 投射
（20.4.9），故可应用（20.6.11）。
